\documentclass{article}

%packacke includes
%-----------------------------------------------------------------------------
\usepackage[utf8]{inputenc}
\usepackage{amsmath}
\usepackage{mathtools}
\usepackage{siunitx}
\usepackage{listings}

% defined macros
%-----------------------------------------------------------------------------
% Ceil and Floor
\DeclarePairedDelimiter\ceil{\lceil}{\rceil}
\DeclarePairedDelimiter\floor{\lfloor}{\rfloor}

\DeclareSIUnit{\bit}{bit}


% Set spellcheck for TexStudio
%!TeX spellcheck = en_US

% define header
%-----------------------------------------------------------------------------
\title{Hello LaTeX!}
\author{Your Name}
\date{\today}

\begin{document}

\maketitle

\newpage

\tableofcontents

\section{Introduction}
Introduction

\newpage
\section{Algorithms}

Here we want to describe the developed algorithms. The idea was to create the union $U$, 
where every processor $r$ has calculated a subset $S_r\subset C$ from a known set $C$ of fixed size. The union
\begin{align}
    U = \bigcup_{r=0}^{p-1} S_r
\end{align}
should be available for every processor. $p$ is here the number of available processes. \\
A straight forward approach to tackle this problem is to create a bitmap for each process, give the elements of $C$ an explicit ordering
and map each element to a bit of the bitmap. An element $e$ is then defined as beeing in the subset $S_r$, if the the corresponding 
bit is set to 1 and set to 0 otherwise. The union can than easily be calculated by performing a bitwise OR-operation over all 
bitmaps from all processes, as the union over all subset $S_r$ is defined as
\begin{align}
    U = \{x \in C | x \in S_0 \lor x \in S_1 \lor \cdots \lor x \in S_{p-1} \}.
\end{align}
This can be done with MPI using the provided Allreduce function. To store $N$ vertices we need $N$ bits in the bitmap.
As a word size of 64-bit integer is quite common on many architectures, we need $\ceil{\frac{N}{64}}$ 64-bit integer
to represent a subset $S_r$. This can be wasteful, if the cardinality of this subset $S_r$ is quite small and most bits of the bitmap are set to 0. We can introduce here a simple sparse format to save memory and therefore communication bandwidth. By using unsigned 32-bit integer, we can represent subsets of $C$ as an explicit list of the selected integer. Here we only can represent $2^{32}$ distinct elements so the cardinality of $C$ is bounded by $2^{32} \approx 4.2\cdot10^9$. The benefit is now, that the size of the sparse representation only depends on the cardinality of the subset $S_r$. We can calculate the memory footprint $M_\mathrm{Sparse}$ of this representation as being 
\begin{align}
	M_\mathrm{Sparse} = m \cdot \qty{32}{\bit},
\end{align}
where $m = |S_r|$ is the cardinality of $S_r$. With the approximate memory footprint of the dense representation $M_\mathrm{Dense} = N \cdot \qty{1}{bit}$, we can estimate, that the sparse variant is more memory efficient, if
\begin{align}
	m \leq \frac{1}{32} = 3.125 \%.
\end{align}
The idea for Algorithm \ref{algo:Algorithm1} is based on the 



\begin{lstlisting}[caption={Algorithm 1}, captionpos=b, label=algo:Algorithm1]
asdsad
asdasd
\end{lstlisting}

\section{Math Test}
Here is a simple equation:
\begin{align*}
    F = G ~ \frac{m M}{r^2}
\end{align*}


\end{document}